\documentclass[10pt,reqno]{amsart}

\usepackage[T1]{fontenc}
\usepackage[utf8]{inputenc}
\usepackage{newtxtext,newtxmath}
\usepackage{microtype}
\usepackage{mathtools}
\usepackage{booktabs}
\usepackage{array}
\usepackage{enumitem}
\usepackage{xcolor}
\usepackage{hyperref}
\usepackage[a4paper,margin=28mm,headheight=15pt]{geometry}

\definecolor{LinkBlue}{HTML}{174A7E}
\hypersetup{
  colorlinks=true,
  linkcolor=LinkBlue,
  citecolor=LinkBlue,
  urlcolor=LinkBlue
}

\setlength{\parskip}{2pt}
\setlist[itemize]{leftmargin=6mm,itemsep=1.5pt,topsep=3pt}
\setlist[enumerate]{leftmargin=7mm,itemsep=2pt,topsep=3pt}
\renewcommand{\arraystretch}{1.15}
\emergencystretch=2em

\newtheorem{theorem}{Theorem}[section]
\newtheorem{proposition}[theorem]{Proposition}
\newtheorem{lemma}[theorem]{Lemma}
\newtheorem{corollary}[theorem]{Corollary}
\theoremstyle{definition}
\newtheorem{definition}[theorem]{Definition}
\newtheorem{remark}[theorem]{Remark}

\title[Six-Point Plateau Obstruction to Degree-Four Hermite Metric Repairs]
{A Six-Point Positive-Weight Obstruction to Degree-Four Hermite Curvature Repairs on Manifolds Without Conjugate Points}
\author{Argus Open Research Project}
\date{August 29, 2026}
\subjclass[2020]{Primary 53C22; Secondary 53C20}
\begin{document}
\begin{abstract}
In the construction of Riemannian metrics without conjugate points whose geodesic rays fail opposite-end divergence, we prove that the degree-four compactly supported Hermite transverse blister repair over $[-16,16]^5$ cannot satisfy the preregistered sectional-curvature gate $K\le 1$: six positive rational plateau weights simultaneously annihilate all five free coefficient directions and yield a strictly negative weighted slack, forcing every real coefficient vector in $\mathbb{R}^5$ to violate $K\le 1$ at one of six exact rational evaluation points.
\end{abstract}
\maketitle
\section{Main Farkas Obstruction}
\begin{theorem}
Let $g_c$ ($c\in\mathbb{R}^5$) be the five-parameter degree-four Hermite compactly supported metric perturbation preserving the full first jet along the central geodesic. There exist six explicit rational evaluation points $q_1,\dots,q_6$ and positive rational weights $w_1,\dots,w_6\in\mathbb{Q}_{>0}$ such that the linearized coefficient matrix $B\in\mathbb{Q}^{6\times 5}$ satisfies $w^{\mathsf T}B = 0$ and the baseline curvature excess vector $b_0$ satisfies $w^{\mathsf T}b_0 > 0$. Consequently, for every $c\in\mathbb{R}^5$, $\max_{1\le i\le 6} K_{g_c}(q_i) > 1$.
\end{theorem}
\end{document}
